\documentclass[conference]{IEEEtran}
\IEEEoverridecommandlockouts
% The preceding line is only needed to identify funding in the first footnote. If that is unneeded, please comment it out.
\usepackage{cite}
\usepackage{amsmath,amssymb,amsfonts}
\usepackage{algorithmic}
\usepackage{textcomp}
\usepackage{xcolor}
\usepackage{silence}
\usepackage{amsthm}
\usepackage{enumitem}
\usepackage{eufrak}
\usepackage{tabularx}
\usepackage{lipsum}
\usepackage{pifont}
\usepackage{yfonts}
\usepackage{booktabs}
\usepackage{float}
\usepackage{amsmath}
\usepackage{amssymb}
\usepackage{comment}
\usepackage{amsfonts}
\usepackage{csquotes}
\usepackage{fancyhdr}
\usepackage{graphicx}
\usepackage{listings}
\usepackage{mathrsfs}
\usepackage{setspace}
\usepackage{textcomp}
\usepackage{multicol}
\usepackage{blindtext}
\usepackage{enumerate}
\usepackage{kantlipsum}
\def\BibTeX{{\rm B\kern-.05em{\sc i\kern-.025em b}\kern-.08em
    T\kern-.1667em\lower.7ex\hbox{E}\kern-.125emX}}
\renewcommand{\refname}{Referencias}
\renewcommand{\lstlistingname}{Listado}
% Marcador de contenido pendiente (se muestra en rojo).
\newcommand{\PEND}[1]{{\color{red}[#1]}}
\begin{document}

\title{Validación de una microrred DC de 48 V emulada: condiciones nominales y ante perturbaciones sobre STM32 y ESP32\\
}

\author{\IEEEauthorblockN{Sofía Parra Rozo}
\IEEEauthorblockA{\textit{Departamento de Ingeniería Eléctrica y Electrónica} \\
\textit{Universidad de los Andes}\\
Bogotá, Colombia \\
s.parrar2@uniandes.edu.co}
\and
\IEEEauthorblockN{Tomás Durán Cardoso}
\IEEEauthorblockA{\textit{Departamento de Ingeniería Eléctrica y Electrónica} \\
\textit{Universidad de los Andes}\\
Bogotá, Colombia\\
t.duranc@uniandes.edu.co}
}

\maketitle

\begin{abstract}
Este documento presenta la validación de una microrred DC de 48\,V emulada sobre microcontroladores, con cuatro subsistemas (arreglo fotovoltaico, turbina eólica, banco de baterías y carga variable) y un nivel supervisorio. Cada subsistema vive en una STM32F303K8 que integra su fuente, un convertidor DC-DC y su control local, y un gateway ESP32-S3 los acopla, ejecuta una máquina de estados de supervisión (recorte proporcional de generación, deslastre de carga por niveles, modo seguro y arranque en negro) y publica por MQTT. Respecto a la entrega anterior se corrigió el seguimiento del punto de máxima potencia del PV, que quedaba en circuito abierto al amanecer, y se rediseñó el deslastre. La validación usa criterios de aceptación cuantitativos y separa de forma explícita el hardware real de la emulación. En una corrida de 3.3 días simulados el PV arrancó solo en 3 de 3 amaneceres, el bus se mantuvo en 48.00\,V y el SOC mínimo fue de 24.5\,\%, frente a 13.0\,\% con el deslastre anterior. En excedente, el recorte proporcional redujo 70\,\% el tiempo con el bus sobre 53\,V y subió el tiempo en banda de 47 a 82\,\%, pero no evitó un pico de $\approx$58\,V con la batería llena, y el arranque en negro se completó en unos 25\,s. Los detalles de implementación están en los anexos técnicos y en el repositorio del proyecto.
\end{abstract}

\begin{IEEEkeywords}
microrred DC, emulación, STM32, ESP32, MQTT, SCADA, MPPT, deslastre de carga
\end{IEEEkeywords}

\section{Introducción}

Las microrredes DC resultan atractivas para comunidades aisladas o con red débil, porque sus fuentes (paneles, baterías, rectificadores de pequeñas turbinas) y buena parte de sus cargas ya son de corriente continua. Antes de invertir en equipos de potencia conviene probar cómo se comportan juntos los generadores, el almacenamiento y la demanda, y cómo los gobierna una estrategia de supervisión.

Este trabajo construye esa plataforma de prueba para un bus DC de 48\,V con cuatro subsistemas: un arreglo fotovoltaico, una turbina eólica con rectificador pasivo, un banco de baterías 10s5p que forma el bus y una carga variable (perfiles residencial e industrial). Los modelos se basan en las investigaciones de los demás subgrupos del curso~\cite{aristizabal2026pv,gonzalez2026eolica,vega2026bateria,hernandez2026carga}. Aquí no se replican los equipos físicos sino que se \emph{emulan}: cada subsistema reside en una STM32F303K8 que integra el modelo de su fuente, un convertidor DC-DC promediado y su lazo de control local, y un gateway ESP32-S3 hace de maestro I2C, ejecuta la supervisión (SCADA) y publica por MQTT hacia un registrador y un tablero en vivo.

El objetivo de esta entrega es validar el sistema integrado, con criterios de aceptación cuantitativos fijados \emph{antes} de medir, en condiciones nominales y ante perturbaciones (excedente, déficit, arranque en negro). Las contribuciones son: (i) una arquitectura de una STM32 por subsistema con una trama común de convertidor, de modo que un convertidor físico pueda reemplazar a la emulación sin cambiar el SCADA; (ii) un SCADA como máquina de estados con recorte proporcional de generación y deslastre de carga por niveles; (iii) la corrección de un fallo del seguimiento del punto de máxima potencia (MPPT) del PV al amanecer; y (iv) una validación experimental que separa de forma explícita lo que es hardware real de lo que es emulación.

Para respetar la extensión, este documento presenta el sistema y los resultados; los modelos con sus ecuaciones y curvas, el convertidor universal, el bus I2C, los lazos de control, las reglas completas del SCADA, el detalle de los experimentos de la entrega anterior y la comparación con las arquitecturas de referencia están en los anexos. El código, los datos y las figuras están en el repositorio del proyecto (Anexo~\ref{app:repo}).

\section{Sistema Emulado}
\label{sec:sistema}

\subsection{Arquitectura}

La Fig.~\ref{fig:arq} resume la arquitectura y la Fig.~\ref{fig:montaje} muestra el montaje y el tablero. Las cuatro STM32 son esclavas I2C (100\,kHz) del ESP32-S3; cada una responde en dos direcciones, una de telemetría de planta y convertidor (0x10--0x13) y otra de control y comandos (0x20--0x23). El ESP32-S3 interroga las placas cada segundo, calcula la potencia neta del resto de la microrred y se la envía a la batería, ejecuta el SCADA y publica todo por MQTT. Un script de Python registra los mensajes crudos en un CSV y alimenta un tablero en vivo.

\begin{figure}[t]
\centering
\includegraphics[width=\linewidth]{arquitectura.png}
\caption{Arquitectura del sistema: cuatro STM32 (planta, convertidor y control local) conectadas por I2C al ESP32-S3 (SCADA), que publica por Wi-Fi/MQTT hacia el PC.}
\label{fig:arq}
\end{figure}

\begin{figure}[t]
\centering
\includegraphics[width=0.48\linewidth]{montaje_real.png}\hfill
\includegraphics[width=0.48\linewidth]{scada_captura.png}
\caption{Montaje real (izquierda: cuatro STM32 y el ESP32-S3) y tablero SCADA en vivo (derecha), en operación normal con $V_{bus}=48.0$\,V.}
\label{fig:montaje}
\end{figure}

Los cuatro subsistemas comparten un reloj de simulación en el que un día de 24\,h dura 360\,s (1\,h $=15$\,s) y arrancan a las 6\,h simuladas. El SOC de la batería se integra con una aceleración fija de $60\times$ respecto al tiempo real. La banda de operación es $48\,\text{V}\pm5\,\%$ (45.6--50.4\,V).

\subsection{Subsistemas, convertidor y control local}

La Tabla~\ref{tab:sub} resume cada subsistema. Todos exponen la misma trama de convertidor (tensión de entrada, tensión y corriente de bus, ciclo de trabajo y banderas), definida una sola vez en el firmware y reutilizada por las cuatro placas, y el convertidor se integra en la misma interrupción de 10\,kHz que el modelo de la fuente y el lazo de control. Las ecuaciones, curvas y circuitos equivalentes están en el Anexo~B, y los lazos y sus parámetros en el Anexo~E.

\begin{table}[t]
\caption{Subsistemas emulados (cada uno en una STM32F303K8)}
\label{tab:sub}
\centering
\scriptsize\setlength{\tabcolsep}{2.5pt}
\begin{tabular}{|l|p{2.0cm}|p{1.5cm}|p{2.0cm}|}
\hline
\textbf{Sub.} & \textbf{Fuente emulada} & \textbf{Convertidor} & \textbf{Control local} \\
\hline
PV & Renogy 100\,W, 1\,diodo, $G(h)$ senoidal & Boost & MPPT P\&O + PI de $V_{pv}$ \\ \hline
Eólica & PMSG + rectificador, Thevenin $V_{oc}(f)$, $R_{eq}(f)$ & Boost & MPPT P\&O + PI (mismo módulo del PV) \\ \hline
Batería & Thevenin 1.er orden, 10s5p, 14.5\,Ah & Bidireccional & PI de $V_{bus}$ + PI de corriente \\ \hline
Carga & Potencia constante, perfiles res./ind. (150\,W máx.) & Emulador de corriente & Lazo abierto con $\lambda\in[0,1]$ \\ \hline
\end{tabular}
\end{table}

\textbf{Modelos principales.} El PV genera $G(h)=1000\sin\!\big(\pi(h-6)/12\big)$ entre las 6 y las 18\,h, con $T_{cell}=T_a+\tfrac{NOCT-20}{800}G$ e $I=I_{sc}\tfrac{G}{1000}\,[1+\alpha(T_{cell}-25)]$; a 1000\,W/m$^2$ el modelo da $\approx$88.6\,W en $V_{mp}\approx18$\,V. La turbina es una fuente de Thevenin $V_{dc}=V_{oc}(f)-R_{eq}(f)\,I$ con $f$ fijada por el viento (arranque 3, nominal 11 y corte 25\,m/s) y un máximo ideal de 88.6\,W a 25.5\,Hz. La batería es un Thevenin de primer orden, $V_{bat}=V_{oc}(\mathrm{SOC})-IR_0(\mathrm{SOC})-V_1$, con
\begin{equation}
\mathrm{SOC}_{k+1}=\mathrm{SOC}_k-\frac{a\,\Delta t\,I}{3600\,Q}\times100,\quad a=60,\ Q=14.5\,\mathrm{Ah},
\end{equation}
y limita su corriente a 7\,A reducida por SOC, tensión y temperatura. La carga es de potencia constante, $P=\lambda\,p_u(h)\,P_{max}$ con $P_{max}=150$\,W, e $I_{dc}=P/(\eta(FP)\,V_{bus})$ con $\eta(FP)=0.96-0.30\,(1-FP)^2$.

\textbf{Corrección del MPPT del PV (S14).} En las corridas de S13 el PV quedó en circuito abierto al amanecer en tres ocasiones. La causa fue el P\&O: la regla de inversión del sentido solo actúa cuando la potencia \emph{baja}, y de noche ($p=0$) nunca se cumple, de modo que la referencia de tensión se pegaba a su límite superior (23.8\,V, por encima de $V_{oc}\approx22.3$\,V) y el panel no arrancaba. La corrección tiene dos partes: si $p<0.5$\,W la referencia se estaciona en su valor inicial y el sentido se fija en descenso, y si la referencia toca un límite se invierte el sentido. Esto se verificó primero en simulación sobre tres días simulados y luego en hardware (Sección~\ref{sec:p0}).

\subsection{SCADA}

El SCADA corre en el ESP32-S3 una vez por segundo. No tiene lazos continuos ni filtros: es una máquina de estados con histéresis y tiempos mínimos que solo manda a cada STM un modo (normal, limitado o seguro) y, cuando corresponde, un tope de potencia o un nivel de carga. Sus estados y reglas S0--S9 se resumen en la Tabla~\ref{tab:estados} y se detallan en el Anexo~F.

\begin{table}[t]
\caption{Estados del SCADA}
\label{tab:estados}
\centering
\scriptsize\setlength{\tabcolsep}{3pt}
\begin{tabular}{|l|c|p{3.6cm}|}
\hline
\textbf{Estado} & \textbf{Regla} & \textbf{Acción} \\
\hline
NORMAL & S0 & Bus en banda; generadores en MPPT \\ \hline
EXCEDENTE & S1--S4 & Topes de potencia a la PV y luego a la eólica \\ \hline
DEFICIT & S5 & Batería descargando y bus bajo (sostenido) \\ \hline
Deslastre & S6, S7 & Carga reducida por niveles según el SOC \\ \hline
SEGURO & S8 & Batería sin respuesta: se liberan los topes \\ \hline
ARRANQUE & S9 & Secuencia de arranque en negro, 4 pasos de 6\,s \\ \hline
\end{tabular}
\end{table}

\textbf{Excedente (S1--S4).} Si el bus supera 50.4\,V durante dos ciclos, S1 entra a EXCEDENTE y fija un tope absoluto a los generadores (primero la PV, hasta un piso de 15\,W, y luego la eólica, con piso de 13.4\,W). En S14 el escalón de recorte de S2 es proporcional al exceso de tensión, $\Delta P=\min\{30,\,5+10\,(V_{bus}-50.4)\}$\,W, y se aplica cada 6\,s, o cada 2\,s si el bus supera 53\,V; en S13 era fijo (5\,W). S3 suelta un escalón cuando la batería vuelve a descargar y S4 devuelve los generadores a MPPT.

\textbf{Déficit y deslastre (S5--S7).} En S13 el deslastre tenía un solo nivel (carga al 60\,\% con SOC $<20$\,\%) y no evitó que el SOC bajara a 13.0\,\%. En S14 es escalonado y anticipado: la carga se lleva al 60, 40 y 25\,\% de su perfil cuando el SOC baja de 30, 25 y 20\,\%, y se restaura un nivel cada vez que el SOC supera 40, 35 y 30\,\% (histéresis de 10 puntos), con un tiempo mínimo de 10\,s entre cambios.

\textbf{Arranque en negro (S9).} Por comando, cuatro pasos de 6\,s: todo en modo seguro (bus caído emulado), la batería forma el bus, la PV y la eólica reconectan en MPPT y por último se reconecta la carga.

\subsection{Qué es hardware y qué es emulación}

La Tabla~\ref{tab:hw} separa lo que es físico de lo que es modelo. \emph{Todo lo que ocurre entre una fuente y el bus es software}: no hay etapa de potencia, y por eso este trabajo valida la supervisión y el control, no la transferencia de potencia.

\begin{table}[t]
\caption{Hardware real frente a emulación}
\label{tab:hw}
\centering
\scriptsize\setlength{\tabcolsep}{3pt}
\begin{tabular}{|p{1.45cm}|p{2.7cm}|p{2.6cm}|}
\hline
 & \textbf{Hardware real} & \textbf{Emulado (firmware)} \\
\hline
Cómputo y comunicación & 4 STM32F303K8, ESP32-S3, bus I2C a 100\,kHz, Wi-Fi/MQTT, PC con registrador y tablero & --- \\ \hline
Fuentes & --- & Irradiancia, viento, perfiles de carga, reloj de 360\,s \\ \hline
Potencia & --- & Convertidores promediados, batería (Thevenin, SOC $\times60$), bus de 48\,V \\ \hline
Perturbaciones & --- & Apagón de E5 (modo seguro); el contactor de aislamiento (E6) no está implementado \\ \hline
\end{tabular}
\end{table}

\section{Metodología Experimental}
\label{sec:metodo}

Las pruebas se hicieron con el sistema integrado y el SCADA activo. El registrador guarda cada mensaje MQTT crudo con su marca de tiempo; todas las métricas se calculan a partir de ese CSV (los datos y los scripts están en el repositorio). Cada experimento es una captura continua con un solo registrador. Los datos de E1 y E2 son de la entrega anterior (S13); P0, E3 y E4 se tomaron con el firmware de S14.

Los criterios de aceptación (Tabla~\ref{tab:crit}) se fijaron antes de repetir los experimentos. Un resultado se marca como cumplido solo si el dato medido lo satisface; el estado se evalúa en la Sección~\ref{sec:resumen}.

\begin{table}[t]
\caption{Experimentos y criterios de aceptación}
\label{tab:crit}
\centering
\scriptsize\setlength{\tabcolsep}{2.5pt}
\begin{tabular}{|c|p{2.0cm}|p{3.5cm}|}
\hline
\textbf{Exp.} & \textbf{Condición} & \textbf{Criterio} \\
\hline
P0 & PV y SCADA v7, carga alternante, $\sim$3 días & 0 de 3 amaneceres con el PV atascado \\ \hline
E1 & Carga res./ind. cada 120\,s & $V_{bus}$ en banda $\ge95\,\%$; SOC mín. $\ge20\,\%$ \\ \hline
E2/E3 & Excedente, sin/con SCADA & $V_{bus,max}\le53$\,V; $\le2$ ciclos S1; SCADA reduce $\ge30\,\%$ el tiempo $>53$\,V \\ \hline
E4 & Carga industrial, SCADA activo & Deslastre en $\le2$\,s tras cruzar el umbral; SOC mín. $\ge20\,\%$ \\ \hline
E5 & Arranque en negro, día y noche & Duración $\le30$\,s; $V_{bus}$ dentro de $\pm10\,\%$ \\ \hline
E6 & Aislamiento de falla por contactor & \PEND{criterio por definir con la rúbrica} \\ \hline
\end{tabular}
\end{table}

\section{Resultados}

\subsection{P0: PV corregido y deslastre escalonado}
\label{sec:p0}

Se registró el sistema completo durante 1194\,s ($\approx$3.3 días simulados) con el firmware corregido del PV y el SCADA con deslastre escalonado, con la carga alternando entre los escenarios residencial e industrial. La Fig.~\ref{fig:p0} resume la corrida.

\begin{figure}[t]
\centering
\includegraphics[width=0.95\linewidth]{fig_P0.png}
\caption{P0: potencia del PV, SOC con los eventos de deslastre (S6) y restauración (S7), y $V_{bus}$. Las líneas punteadas marcan el amanecer simulado.}
\label{fig:p0}
\end{figure}

\textbf{PV.} El PV arrancó solo en los tres amaneceres consecutivos ($t\approx349$, 710 y 1069\,s): ya entregaba potencia en la primera muestra del día (3.5--6.1\,W) y alcanzaba el 10\,\% de su pico en 2.8, 3.4 y 4.9\,s. La referencia de tensión se mantuvo entre 15.8 y 19.4\,V (en 19.4\,V de noche) y nunca se acercó al límite de 23.8\,V que causaba el fallo. El pico diario fue de 88.6\,W en los cuatro días.

\textbf{Deslastre.} El SOC mínimo fue de 24.5\,\% ($t=1095$\,s) y nunca estuvo por debajo de 20\,\%, frente al 13.0\,\% de E4 y al 16.8\,\% de E1 en S13. El nivel 1 (60\,\%) se activó con el SOC en 29.9, 29.9 y 29.8\,\%, el nivel 2 (40\,\%) en 24.9\,\%, y la restauración llegó con el SOC en 40.3 y 40.1\,\% (al 100\,\%) y 35.1\,\% (al 60\,\%), como indica el diseño.

\textbf{Bus y balance.} El bus se mantuvo en 48.00\,V durante toda la corrida (100\,\% dentro de la banda) y la batería estuvo descargando el 56.4\,\% del tiempo. La generación fue del 93.5, 90.3 y 95.8\,\% del consumo en cada uno de los tres días completos, de modo que el SOC tiende a bajar de un día a otro (mínimos diarios de 38.5, 29.6, 26.7 y 24.5\,\%): el deslastre evita agotar la batería, pero no corrige que el sistema esté subdimensionado para este perfil de carga. La corrida no tuvo recortes de generación, por lo que no evalúa el excedente.

\subsection{E1--E5}

E1, E2 y E5 se corrieron en S13 con el firmware anterior (P\&O sin parqueo nocturno y deslastre de un solo nivel) y su detalle está en el Anexo~G; E3 y E4 se repitieron con el firmware de S14. Los resultados de S13 fijan la línea base frente a la cual se mide la mejora.

\textbf{E1, operación nominal.} Con la carga alternando cada 120\,s durante 1010\,s, el bus estuvo en banda el 100\,\% del tiempo (48.00\,V, desviación estándar de 0.00\,V), la batería descargó el 64\,\% del tiempo y el SOC cayó de 48\,\% a un mínimo de 16.8\,\% (Tabla~\ref{tab:e1r}, Fig.~\ref{fig:e1}). La diferencia entre escenarios se explica por el balance de potencia: en el residencial la generación media (65.4\,W) supera a la carga (56.9\,W) y la batería se carga en promedio, mientras que en el industrial la carga (109.4\,W) supera a la generación (80.7\,W) y la batería cubre el déficit el 71\,\% del tiempo. El sistema estuvo en DEFICIT el 58\,\% de las muestras y no pasó por EXCEDENTE, de modo que E1 valida las reglas de déficit pero no las de excedente.

\begin{table}[t]
\caption{E1: métricas por escenario de carga (S13, 1010\,s)}
\label{tab:e1r}
\centering
\scriptsize
\begin{tabular}{|l|c|c|c|}
\hline
\textbf{Métrica} & \textbf{Resid.} & \textbf{Indust.} & \textbf{Global} \\ \hline
$P$ carga media (W) & 56.9 & 109.4 & 81.9 \\ \hline
$P$ generación media (W) & 65.4 & 80.7 & 72.8 \\ \hline
Batería descargando (\%) & 58 & 71 & 64 \\ \hline
$V_{bus}$ en banda (\%) & 100 & 100 & 100 \\ \hline
SOC mínimo (\%) & 16.8 & 18.7 & 16.8 \\ \hline
\end{tabular}
\end{table}

\begin{figure}[t]
\centering
\includegraphics[width=0.85\linewidth]{E1_serie_completa.png}
\caption{E1: tensión de bus, potencias y SOC durante 1010\,s (sombreado: carga industrial).}
\label{fig:e1}
\end{figure}

\textbf{E2 y E3, excedente.} Con la batería llena, el bus se mantuvo en 48.00\,V hasta que su protección por SOC dejó de absorber potencia; sin ella regulando, el sobrante elevó el bus por encima de 50.4\,V. En E2 (S13) los dos episodios llegaron a 57.0 y 55.1\,V y el bus estuvo 9.2\,\% del tiempo fuera de banda. E3 se repitió con el firmware de S14 (recorte proporcional, PV corregido): una corrida de 1009\,s sin SCADA (A) y otra de 1081\,s con SCADA (B), cada una de unos tres días simulados (Tabla~\ref{tab:e3r}, Fig.~\ref{fig:e3}). Las dos corridas no arrancan en la misma fase del día, por lo que se comparan fracciones de tiempo y no instantes; el excedente estuvo presente una fracción similar de cada una (52.9\,\% en A y 53.5\,\% en B).

\begin{table}[t]
\caption{E3 con el firmware de S14: sin SCADA (A) y con SCADA (B)}
\label{tab:e3r}
\centering
\scriptsize
\begin{tabular}{|l|c|c|}
\hline
\textbf{Métrica} & \textbf{A (sin SCADA)} & \textbf{B (con SCADA)} \\
\hline
Tiempo con $V_{bus}>53$\,V & 529\,s (52.4\,\%) & 168\,s (15.5\,\%) \\ \hline
Tiempo dentro de banda & 47.1\,\% & 82.0\,\% \\ \hline
$V_{bus}$ máximo (V) & 58.00 & 57.91 \\ \hline
Recorte máximo (W) & 0 & 90 \\ \hline
Entradas a EXCEDENTE (S1) & --- & 2 \\ \hline
Latencia de S1 tras $V_{bus}>50.4$\,V & --- & 0.7--1.1\,s \\ \hline
\end{tabular}
\end{table}

\begin{figure}[t]
\centering
\includegraphics[width=0.92\linewidth]{fig_E3.png}
\caption{E3 (S14): $V_{bus}$ sin SCADA (A) y con SCADA (B), recorte aplicado por el SCADA y potencias. Sombreado: SCADA en EXCEDENTE; banda verde: 45.6--50.4\,V.}
\label{fig:e3}
\end{figure}

El SCADA redujo de 52.4 a 15.5\,\% la fracción de tiempo con el bus sobre 53\,V ($-70\,\%$), subió el tiempo en banda de 47.1 a 82.0\,\% y entró a S1 en $\approx$1\,s, con solo 2 entradas en tres días. No evitó el pico: el bus llegó a 57.91\,V, igual que sin SCADA. Esos picos ocurren cuando la batería está llena ($\mathrm{SOC}\ge99.8\,\%$ y $P_{bat}\approx0$) y el recorte ya llegó a su piso (topes de 15\,W en el PV y 13\,W en la eólica, unos 28--34\,W de generación mínima) que supera la carga de 16--32\,W; ahí el SCADA repite escalones de $+0$\,W sin poder bajar más. Es un límite de dimensionamiento del recorte y de la emulación, que no disipa excedente, y no de la lógica de supervisión.

\textbf{E4, carga industrial.} Con el SCADA activo y SOC inicial de 46.5\,\%, la batería subió a $\approx$100\,\% en unos 500\,s por la generación, el SOC mínimo fue de 45.6\,\% y el deslastre (S6) no se activó (Fig.~\ref{fig:e4}). Por eso esta corrida \emph{no permite evaluar} la latencia del deslastre ni el SOC mínimo bajo estrés; esos criterios se evalúan con P0, donde S6 actuó a 29.9\,\% y el SOC mínimo fue de 24.5\,\% (13.0\,\% con el deslastre de S13). E4 sirvió como segunda corrida de excedente con carga industrial: el SCADA recortó hasta 118\,W, el bus estuvo en banda el 90.4\,\% del tiempo, sobre 53\,V el 5.6\,\% (57\,s) y con un máximo de 57.11\,V, y hubo 2 entradas a S1 con latencia de 1\,s.

\begin{figure}[t]
\centering
\includegraphics[width=0.92\linewidth]{fig_E4.png}
\caption{E4 (S14): tensión del bus, potencias y SOC con carga industrial. El SOC nunca se acercó al umbral de deslastre.}
\label{fig:e4}
\end{figure}

\textbf{E5, arranque en negro.} La secuencia se completó en 24.8\,s de día y 25.4\,s de noche, con pasos de entre 5.6 y 7.4\,s y sin eventos S8. Con la batería limitada a $\pm2$\,A en modo seguro, el bus subió a 52.92\,V de día ($+4.9$\,V, durante unos 4\,s) y bajó a 45.16\,V de noche ($-2.9$\,V, unos 2\,s), y en ambos casos regresó a 48.00\,V sin intervención. De día entraban 107\,W de excedente, más de lo que la batería limitada puede absorber, y de noche la carga (93\,W) superaba lo que puede entregar; la carga bajó a su perfil mínimo (10--15\,W) durante unos 18\,s y volvió a su perfil entre 1.1 y 1.7\,s después del último paso (Fig.~\ref{fig:e5}). El apagón es emulado: E5 valida la secuencia de reconexión y no un arranque desde un bus realmente desenergizado.

\begin{figure}[t]
\centering
\includegraphics[width=\linewidth]{E5_blackstart.png}
\caption{E5: arranque en negro de día (izquierda) y de noche (derecha). Sombreado: secuencia de 4 pasos; banda: 45.6--50.4\,V.}
\label{fig:e5}
\end{figure}

\subsection{Resumen frente a los criterios}
\label{sec:resumen}

\begin{table}[t]
\caption{Resultados frente a los criterios de aceptación}
\label{tab:res}
\centering
\scriptsize\setlength{\tabcolsep}{2.5pt}
\begin{tabular}{|c|p{3.9cm}|c|}
\hline
\textbf{Exp.} & \textbf{Resultado medido} & \textbf{Estado} \\
\hline
P0 & 0 de 3 amaneceres atascados & \ding{51} \\ \hline
E1 & $V_{bus}$ en banda 100\,\%; SOC mín. 16.8\,\% (S13); 24.5\,\% en P0 (S14) & \ding{51}\,/\,\ding{55}\,/\,\ding{51} \\ \hline
E2/E3 & $V_{bus,max}$ 57.91\,V; reducción del tiempo $>53$\,V: 70\,\%; 2 entradas a S1 & \ding{55}\,/\,\ding{51}\,/\,\ding{51} \\ \hline
E4 & SOC mín. 45.6\,\%, S6 no se activó; en P0, S6 a 29.9\,\% y SOC mín. 24.5\,\% & no evaluable en E4; \ding{51} en P0 \\ \hline
E5 & 24.8 / 25.4\,s; $+10.3\,\%$ de día, $-5.9\,\%$ de noche & \ding{51}\,/\,$\approx$ \\ \hline
E6 & No realizado & \PEND{pendiente} \\ \hline
\end{tabular}
\end{table}

El comportamiento supervisorio queda demostrado en lo que se midió: el PV ya no se atasca al amanecer, el bus permanece en banda mientras la batería lo regula, y el deslastre escalonado mantuvo el SOC por encima de 20\,\% en P0. Con los datos de S14, el SCADA cumple la reducción del tiempo sobre 53\,V (70\,\% frente al 30\,\% exigido) y el límite de ciclos S1, pero no el tope de $V_{bus}\le53$\,V (57.91\,V), por el piso del recorte con la batería llena. E4 no ejercitó el deslastre, que se evalúa con P0. Con los datos de S13, el exceso de $+10.3\,\%$ del bus en el arranque en negro diurno supera por muy poco el $\pm10\,\%$.

\section{Comparación con las Arquitecturas de Referencia}

Lo reportado es una emulación de señal y control sobre microcontroladores, no una réplica de los bancos físicos de los demás subgrupos. El PV equivale únicamente al modelo simplificado de la referencia~\cite{aristizabal2026pv}, con irradiancia sintética. La turbina simplifica la regulación del bus que el banco real~\cite{gonzalez2026eolica} no logra sin una etapa activa. La batería es el subsistema más alineado con su referencia~\cite{vega2026bateria} (14.5\,Ah, límites de 0.5C y resistencia creciente con la descarga). La brecha mayor es la electrónica de potencia~\cite{pinto2026peb}: aquí no hay transferencia de potencia física. El SCADA reemplaza el PLC con Modbus TCP de la referencia~\cite{alvarez2026scada} por un gateway ESP32 con MQTT, una alternativa más liviana, no equivalente. La comparación completa está en el Anexo~H.

\section{Conclusiones}

Se construyó una microrred DC de 48\,V emulada en la que cuatro subsistemas con modelo, convertidor y control local corren sobre un mismo reloj, con un SCADA que los acopla y supervisa. Se demostró que (i) la corrección del P\&O elimina el atasco del PV al amanecer (3 de 3 amaneceres sin fallo), (ii) el bus se mantiene en 48.00\,V mientras la batería lo regula y (iii) el deslastre por niveles sube el SOC mínimo de 13.0 a 24.5\,\%. En excedente, el recorte proporcional de S14 redujo 70\,\% el tiempo con el bus sobre 53\,V y subió el tiempo en banda de 47 a 82\,\%, pero no evitó el pico de $\approx$58\,V cuando la batería está llena y el recorte llega a su piso; el arranque en negro se completa en unos 25\,s.

\textbf{Limitaciones.} Cada experimento es una sola corrida; no hay repeticiones para estimar dispersión. El apagón de E5 y todas las fuentes y convertidores son emulados, y el excedente no se disipa en el modelo. La microrred está subdimensionada para el perfil casa--industria (generación del 90--96\,\% del consumo). La corrida E4 de S14 no llevó el SOC al umbral de deslastre, y el aislamiento de falla por contactor (E6) está pendiente.

\textbf{Trabajo futuro.} Repetir E4 partiendo de un SOC cercano al umbral de deslastre; añadir una carga de disipación o un piso de recorte menor para acotar el pico de tensión; implementar el aislamiento de falla por contactor; medir los lazos de la turbina, la batería y la carga en hardware, y sustituir progresivamente los subsistemas emulados por equipos físicos aprovechando la trama común del convertidor, tras resolver las protecciones, la puesta a tierra y el dimensionamiento.

\begin{thebibliography}{00}

\bibitem{gonzalez2026eolica}
J. D. González Cubides, J. D. Roa Ballén, and C. D. González Mayorga, ``Emulación experimental y supervisión SCADA de un subsistema eólico de 300 W para microrredes DC basado en PMSG,'' Informe técnico, Departamento de Ingeniería Eléctrica y Electrónica, Universidad de los Andes, Bogotá, Colombia, 2026.

\bibitem{aristizabal2026pv}
L. Á. Aristizábal Correa, J. F. Angulo Rubiano, and S. J. Fonseca Boada, ``Emulación de Arreglos Fotovoltaicos para Microrredes DC de Laboratorio: Pipeline en Python, Comparación de Modelos Eléctricos e Integración SCPI--Modbus TCP,'' Informe técnico, Departamento de Ingeniería Eléctrica y Electrónica, Universidad de los Andes, Bogotá, Colombia, 2026.

\bibitem{alvarez2026scada}
M. Alvarez Valenzuela, A. F. Mosquera Gómez, and Ó. M. Cristancho Florez, ``Informe Final: SCADA,'' Taller de Potencia, Departamento de Ingeniería Eléctrica y Electrónica, Universidad de los Andes, Bogotá, Colombia, 2026.

\bibitem{hernandez2026carga}
J. F. Hernández, D. E. Niño Jimenez, and J. J. Castillo Urbano, ``Metodología de Emulación de Perfiles de Demanda en Microrredes DC mediante Integración C-HIL y Cargas Programables,'' Informe técnico, Departamento de Ingeniería Eléctrica y Electrónica, Universidad de los Andes, Bogotá, Colombia, 2026.

\bibitem{pinto2026peb}
J. D. Pinto Cardenas and J. J. Vasquez Polaina, ``Informe final: Bloque Modular de Electrónica de Potencia (PEBB) basado en un Convertidor DC--DC en Configuración de Medio Puente con IGBT para Microrredes DC,'' Departamento de Ingeniería Eléctrica y Electrónica, Universidad de los Andes, Bogotá, Colombia, 2026.

\bibitem{vega2026bateria}
L. D. Vega Moreno, J. S. Torres Soler, and J. J. Carrera Luna, ``Integración y caracterización de un sistema de almacenamiento 10s5p basado en Ion-Litio para microrredes DC,'' Departamento de Ingeniería Eléctrica y Electrónica, Universidad de los Andes, Bogotá, Colombia, 2026.

\bibitem{bader2019pv}
S.~Bader, X.-Y.~Ma, and B.~Oelmann, ``One-diode photovoltaic model parameters at indoor illumination levels -- A comparison,'' \emph{Solar Energy}, vol.~180, pp.~707--716, 2019. DOI: 10.1016/j.solener.2019.01.048.

\bibitem{han2024battery}
S.~Han, ``Optimal Charging Current Protocol with Multi-Stage Constant Current Using Dandelion Optimizer for Time-Domain Modeled Lithium-Ion Batteries,'' \emph{Applied Sciences}, vol.~14, no.~23, p.~11320, 2024. DOI: 10.3390/app142311320.

\bibitem{benavides2025eolica}
D.~E.~Benavides-Mendoza, C.~L.~Trujillo-Rodríguez, and Ó.~D.~Montoya-Giraldo, ``Modelado, diseño y control de un sistema de generación eólica basado en MPPT con convertidor electrónico de potencia para la integración de una turbina a una microrred,'' \emph{TecnoLógicas}, vol.~28, no.~62, pp.~1--28, 2025. DOI: 10.22430/22565337.3130.

\bibitem{rodriguezlicea2018cpl}
M.-A.~Rodríguez-Licea, F.-J.~Pérez-Pinal, J.-C.~Nuñez-Perez, and C.-A.~Herrera-Ramirez, ``Nonlinear Robust Control for Low Voltage Direct-Current Residential Microgrids with Constant Power Loads,'' \emph{Energies}, vol.~11, no.~5, p.~1130, 2018. DOI: 10.3390/en11051130.

\bibitem{ontibon2025pebb}
J.~Ontibón \emph{et al.}, ``Power Electronics Building Block for Power-Hardware-In-The-Loop and Its Application in Microgrid Research,'' 2025.

\bibitem{anthropic2026claude}
Anthropic, ``Claude (Sonnet),'' modelo de lenguaje de gran escala utilizado para la revisión de formato del documento, verificación de consistencia de citas, y revisión de gramática y ortografía, Anthropic, 2026. [En línea]. Disponible: https://www.anthropic.com/claude

\end{thebibliography}


\appendices
\section{Repositorio y anexos técnicos}
\label{app:repo}
Código, datos y anexos técnicos están en \texttt{https://github.com/soofffp/VITA-DC-Microgrid} (etiqueta \texttt{entrega-S14}). El documento \texttt{docs/anexos\_tecnicos.pdf} reúne los modelos y curvas (Anexo~B), el convertidor y el bus I2C (C, D), los lazos de control (E), las reglas del SCADA (F), el detalle de los experimentos de S13 (G) y la comparación con las referencias (H). El \texttt{README.md} explica cómo compilar y cargar cada placa y cómo correr el registrador y el análisis.

\end{document}
